泛函行列式的固有时表示把一圈紫外问题转化为短固有时的局域渐近展开。对欧氏正谱二次算符 \(\Delta\)，先选取逆长度减除尺度 \(\mu_\kappa\,[\mathrm{m^{-1}}]\)。谱定理把标量恒等式 \(\ln(\lambda/\mu_\kappa^2)\) 提升为


\begin{proposition}[正算子行列式的固有时表示]
设欧氏二次算子为正自伴，先以有限模或热核截断定义相对行列式，再按正文做紫外减除；非正谱需另选解析延拓。则有
\begin{equation}
 {
 \Tr\ln\frac{\Delta}{\mu_\kappa^2}
 =-\int_0^\infty\frac{\dd s}{s}
 \Tr\!\left(\ee^{-s\Delta}-\ee^{-s\mu_\kappa^2}\right).}
\label{eq:proper-time-functional-dp11}
\end{equation}
\end{proposition}
\begin{proof}
正文\eqnrefs{eq:one_loop_effective_action_general} 表明一圈修正由二次涨落算符的 $\Tr\ln$ 决定。为把紫外问题改写成可按短距离展开的积分，先 Wick 转动到正谱欧氏算符 $\Delta$，再把标量对数恒等式逐个作用在 $\Delta$ 的本征值上。

对正数 $\lambda>0$ 定义
\begin{equation*}
 F(\lambda)
 =-\int_0^\infty\frac{\dd s}{s}
 \left(\ee^{-s\lambda}-\ee^{-s\mu_\kappa^2}\right).
\end{equation*}
对 $\lambda$ 求导，
\begin{equation*}
 \frac{\dd F}{\dd\lambda}
 =\int_0^\infty\dd s\,\ee^{-s\lambda}
 =\frac1\lambda,
\end{equation*}
并且 $F(\mu_\kappa^2)=0$，所以
\begin{equation*}
 F(\lambda)=\ln\frac{\lambda}{\mu_\kappa^2}.
\end{equation*}
若 $\Delta|n\rangle=\lambda_n|n\rangle$，谱定理给出
\begin{align*}
 \Tr\ln\frac{\Delta}{\mu_\kappa^2}
 &=\sum_n\ln\frac{\lambda_n}{\mu_\kappa^2}\\
 &=-\sum_n\int_0^\infty\frac{\dd s}{s}
 \left(\ee^{-s\lambda_n}-\ee^{-s\mu_\kappa^2}\right)\\
 &=-\int_0^\infty\frac{\dd s}{s}
 \Tr\!\left(\ee^{-s\Delta}-\ee^{-s\mu_\kappa^2}\right),
\end{align*}
这就是正文\eqnrefs{eq:proper-time-functional-dp11}。
\end{proof}

因为 \(\Delta\) 的本征值具有 \(\mathrm{m^{-2}}\) 量纲，固有时 \(s\) 具有 \(\mathrm{m^2}\) 量纲；大本征值对应 \(s\to0^+\)，所以紫外发散完全由热核的短时系数决定。若改用物理动量尺度，则 \(\mu_p=\hbar\mu_\kappa\)，不能把二者直接混入同一个无量纲比值。

在平直欧氏空间中，对 Laplace 型算符定义
\begin{equation}
 {
 \Delta=-\left(\nabla^2+E\right),
 \qquad
 \nabla_\mu=\partial_\mu+\omega_\mu,
 \qquad
 \Omega_{\mu\nu}=[\nabla_\mu,\nabla_\nu].}
 \label{eq:laplace-type-operator-dp11}
\end{equation}
其中 \(E\) 是纤维内部空间上的局域端同态，\(\omega_\mu\) 是联络，\(\Omega_{\mu\nu}\) 是其曲率。对角热核的短时展开为


\begin{proposition}[Laplace 型算子的局部热核系数]
设底空间为平直欧氏空间，算子为光滑 Laplace 型算子，采用正文的联络曲率与势项符号；在远离边界的内点，局部展开在 $s\downarrow0$ 时按渐近意义理解。下式成立；一般曲底空间还须加入 Riemann 曲率项：
\begin{equation}
 {
 K(x,x;s)
 =\frac1{(4\pi s)^{d/2}}
 \left[
 I+sE+s^2\left(
 \frac12E^2+\frac1{12}\Omega_{\mu\nu}\Omega^{\mu\nu}
 +\frac16\nabla^2E
 \right)+\order(s^3)
 \right].}
\label{eq:heat-kernel-a4-flat-dp11}
\end{equation}
\end{proposition}
\begin{proof}
在基点 $x$ 选择 Fock--Schwinger 规范

\begin{equation*}
 (y-x)^\mu\omega_\mu(y)=0,
\end{equation*}
则联络在 $x$ 附近为
\begin{equation*}
 \omega_\mu(y)
 =\frac12(y-x)^\nu\Omega_{\nu\mu}(x)
 +O((y-x)^2).
\end{equation*}
自由热核是
\begin{equation*}
 K_0(x,y;s)=\frac1{(4\pi s)^{d/2}}
 \exp\!\left[-\frac{(x-y)^2}{4s}\right].
\end{equation*}
令 $\xi=y-x$，其高斯矩为
\begin{align*}
 \langle\xi^\mu\xi^\nu\rangle_s
 &=2s\,\delta^{\mu\nu},\\
 \langle\xi^\mu\xi^\nu\xi^\rho\xi^\sigma\rangle_s
 &=4s^2(\delta^{\mu\nu}\delta^{\rho\sigma}
 +\delta^{\mu\rho}\delta^{\nu\sigma}
 +\delta^{\mu\sigma}\delta^{\nu\rho}).
\end{align*}
把 $E$ 与联络项作为固有时有序展开中的插入：一次 $E$ 插入给 $sE$，两次 $E$ 插入给 $s^2E^2/2$；两个线性联络插入结合四阶高斯矩给
\begin{equation*}
 \frac{s^2}{12}\Omega_{\mu\nu}\Omega^{\mu\nu},
\end{equation*}
而 $E(y)$ 的二阶协变 Taylor 项结合二阶高斯矩给
\begin{equation*}
 \frac{s^2}{6}\nabla^2E.
\end{equation*}
将这些项与自由热核前因子合并，就得到正文\eqnrefs{eq:heat-kernel-a4-flat-dp11}。
\end{proof}

无边界体积分中 \(\nabla^2E\) 是全导数，因此四维对数紫外极点真正依赖的局域不变量是 \(E^2\) 与 \(\Omega_{\mu\nu}\Omega^{\mu\nu}\)。

在 \(d=4-2\epsilon\) 中，\eqnrefs{eq:heat-kernel-a4-flat-dp11} 的 \(s^2\) 项与 \(\dd s/s\) 组合成 \(\dd s\,s^{-1+\epsilon}\)，从而直接产生 \(1/\epsilon\) 极点。对实玻色场的一圈高斯行列式，发散部分为
\begin{equation}
 {
 \left.\frac12\Tr\ln\Delta\right|_{\rm div}
 =-\frac1{2(4\pi)^2\bar\epsilon}
 \int\dd^4 x\,
 \tr\!\left[
 \frac12E^2+\frac1{12}\Omega_{\mu\nu}\Omega^{\mu\nu}
 \right].}
 \label{eq:heat-kernel-divergence-master-dp11}
\end{equation}
反交换标量变量和 Dirac 费米变量会因统计性质改变泛函行列式的整体符号与简并度，但允许出现的局域算符基仍由同一短时系数决定。对非阿贝尔规范场，规范场本身与规范固定 Jacobian 产生的 Grassmann 辅助场分别贡献各自的有限维指标迹；这些辅助变量的出现来自规范轨道换元的 Jacobian，而不是额外的物理渐近自由度。
